\section{Conclusion}
We propose \name{}, a framework for disentangling knowledge and application abilities in LLM evaluation through knowledge-matched contrasts.
Our analysis of existing Korean benchmarks reveals a predominance of recall-only questions and the difficulty of isolating application demand.
Through anchor-based generation and iterative validation, \name{} constructs controlled contrasts that vary application demand while preserving required knowledge. 
Evaluations across diverse LLMs reveal increasing application failures that persist even when the required knowledge is provided. 
Collectively, our findings highlight the value of knowledge-matched contrasts for diagnosing model limitations that conventional benchmark scores may obscure.


\newpage
\section*{Limitations} 
While \name{} enables more controlled evaluation of knowledge and application abilities, several limitations remain.

First, our framework currently focuses on multiple-choice questions, which facilitate automated answer validation and evaluation.
However, this restricts its applicability to open-ended tasks, where models must generate answers without predefined options.
Future work could extend \name{} to free-form question answering by incorporating semantic answer validation and evaluation methods that accommodate diverse valid responses.

Second, our proxy-model difficulty screening applies the same acceptance criterion to A2 and A3 candidates, limiting fine-grained control over question difficulty.
Moreover, because candidates are selected based on the performance of a specific proxy model, the retained questions may reflect model-specific weaknesses.
Although evaluations across diverse LLM families demonstrate diagnostic utility beyond the screening model, potential selection bias remains. 
Future work could introduce level-specific thresholds based on repeated proxy evaluations, such as requiring different minimum numbers of incorrect responses across multiple runs, to enable more flexible difficulty calibration.
Additionally, incorporating multiple proxy models or alternative screening configurations could reduce dependence on individual model capabilities while preserving the intended application requirements. 

Lastly, our experiments focus on questions derived from three existing Korean benchmarks.
While these benchmarks span diverse knowledge domains, the generalizability of \name{} to other languages and domains remains to be established.
Future research could extend the framework to a broader range of benchmarks to assess its applicability across different evaluation settings.

\section*{Broader Impact and Ethical Implications}
\name{} can contribute to more interpretable LLM evaluation on Korean benchmarks by distinguishing limitations in knowledge access from those in knowledge application.
Such diagnostic insights may support more targeted model development and reduce reliance on aggregate accuracy as the sole measure of model capability.

However, LLM-generated questions and automated judgments may inherit biases or inaccuracies from source benchmarks and underlying models, potentially affecting the fairness of evaluation.
Therefore, results from \name{} should complement broader evaluations rather than serve as definitive assessments of model capabilities, particularly in high-stakes settings.
In addition, since our framework constructs questions from existing benchmarks, the redistribution of generated materials should respect the licensing and attribution requirements of the original datasets.
We encourage transparent reporting of construction procedures and validation outcomes to support responsible use of the resulting benchmark.

\section*{Acknowledgments}

\noindent\textbf{Dataset Contributors.}
The KLEAR benchmark was validated by the following annotators at DATUMO INC., who assessed each generated question against the human rubric and provided the feedback used to revise it:

\noindent
Hosik Lee, Minseung Song, Jinsol Lee, Chanho Lee, Dahun Han, Nakyung Kim, Nawon Kim, Dongchan Lee, Dahye Kim, Hyeonah Ju, and Younghyun Jo.